

\chapter{Polarization and Restitution:}

\section{Invariant functions, quotients, and finite generation}

\subsection{Invariant functions.}
Let $k$ be an algebraically closed field and let $G$ be an algebraic group acting on an affine variety
$X = \Specm A$, where $A$ is a finitely generated $k$-algebra.
The action of $G$ on $X$ is equivalently given by a rational action of $G$ on $A$ by $k$-algebra automorphisms.
We define the ring of \emph{invariant functions} by
\[
A^G := \{ f \in A \mid g \cdot f = f \text{ for all } g \in G \}.
\]

\begin{remark}[Geometric meaning]
	An invariant function is constant along $G$-orbits.
	Thus $A^G$ measures exactly the functions on $X$ that cannot distinguish points inside a single orbit.
	The philosophy of GIT is that the quotient space should have coordinate ring $A^G$.
\end{remark}

\subsection{Spec vs.\ Specm.}
There are two closely related spaces associated to a ring $A^G$:
\[
\Spec A^G \quad \text{and} \quad \Specm A^G.
\]
The maximal spectrum $\Specm A^G$ corresponds to closed points, while $\Spec A^G$ includes also
generic points encoding specialization relations.

\begin{remark}
	In algebraic geometry, $\Spec A^G$ is the correct object for forming quotients.
	Even when one is interested in closed orbits (points), non-closed orbits play a crucial role, and
	their closures are naturally encoded in $\Spec A^G$, not in $\Specm A^G$.
\end{remark}

\begin{example}
Let $G = \mathbb{Z}/2$ act on $A = k[x,y]$ by
\[
\sigma(x) = y, \qquad \sigma(y) = x.
\]
Then the invariant ring is
\[
A^G = k[x+y,\, xy].
\]


	The morphism
	\[
	\pi : \mathbb{A}^2 \to \mathbb{A}^2, \qquad (x,y) \mapsto (x+y, xy)
	\]
	is $G$-invariant and induces an isomorphism
	\[
	\Spec k[x+y,xy] \cong \Spec A^G.
	\]
	Geometrically, $\pi$ identifies the points $(x,y)$ and $(y,x)$, and separates distinct orbits except
	when $(x,y)$ and $(y,x)$ lie in the same orbit closure.
	In particular, invariant functions separate \emph{closed} orbits but may fail to separate non-closed ones.
\end{example}

\begin{remark}[Key phenomenon]
	The fibers of $\pi$ are precisely orbit closures, not individual orbits.
	This is the fundamental reason the GIT quotient is a \emph{categorical} quotient rather than a geometric one.
\end{remark}

\subsection{Affine schemes and schemes.}
An \emph{affine scheme} is a space of the form $\Spec A$ for a ring $A$.
A \emph{scheme} is obtained by gluing affine schemes along open subsets.
Thus schemes are locally affine objects.

\begin{remark}
	For an affine scheme $X = \Spec A$, one has a canonical identification
	\[
	A \cong \mathcal{O}_X(X).
	\]
	This property characterizes affine schemes among all schemes.
\end{remark}

\subsection{Quotients and finite generation.}
Given an affine $G$-variety $X = \Spec A$, a natural candidate for the quotient is
\[
X /\!/ G := \Spec A^G.
\]
This construction only makes sense if $A^G$ is finitely generated as a $k$-algebra.

\begin{theorem}[Hilbert, classical form]
	If $G$ is a reductive algebraic group acting rationally on a finitely generated $k$-algebra $A$,
	then $A^G$ is finitely generated.
\end{theorem}

\begin{remark}[Why reductivity matters]
	Reductive groups have enough invariant linear functionals to control the growth of invariants.
	For non-reductive groups, invariant rings can be pathologically large.
\end{remark}

\subsection{Hilbert's 14th problem and Nagata's counterexample.}
Hilbert asked whether $A^G$ is finitely generated for \emph{every} linear action of an algebraic group.
This is false.

\begin{theorem}[Nagata]
	There exists a linear action of a unipotent algebraic group $G$ on a polynomial ring $A$
	such that the invariant ring $A^G$ is not finitely generated.
\end{theorem}

\begin{remark}[Conceptual takeaway]
	Finite generation is not a formal property of group actions; it is a deep structural fact tied to
	reductivity.
	GIT works precisely because it restricts attention to reductive group actions, where quotients
	exist as affine schemes of finite type.
\end{remark}
